\documentclass[12pt,oneside]{book}

% ============================================================
% PAQUETES
% ============================================================
\usepackage[utf8]{inputenc}
\usepackage[T1]{fontenc}
\usepackage[spanish,es-nodecimaldot]{babel}

\usepackage[
    top=2.5cm,
    bottom=2.5cm,
    left=3cm,
    right=2.5cm,
    headheight=15pt
]{geometry}

\usepackage{amsmath}
\usepackage{amssymb}
\usepackage{mathtools}

\usepackage{graphicx}
\usepackage{float}
\usepackage{caption}
\usepackage{subcaption}

\usepackage{booktabs}
\usepackage{array}
\usepackage{longtable}
\usepackage{tabularx}

\usepackage{enumitem}

\usepackage{xcolor}
\usepackage[most]{tcolorbox}

\usepackage{fancyhdr}
\usepackage{ragged2e}

\graphicspath{
    {./img/}
    {./graficos/}
    {./figuras/}
}

% ============================================================
% CONFIGURACIÓN GENERAL
% ============================================================
\setcounter{tocdepth}{2}

\setlength{\parindent}{0pt}
\setlength{\parskip}{0.5em}

\setlist[itemize]{
    topsep=0.3em,
    itemsep=0.2em
}
\setlist[enumerate]{
    topsep=0.3em,
    itemsep=0.2em
}

\clubpenalty=10000
\widowpenalty=10000

% ============================================================
% COMANDOS MATEMÁTICOS
% ============================================================
\newcommand{\abs}[1]{\left|#1\right|}
\newcommand{\norm}[1]{\left\lVert#1\right\rVert}
\newcommand{\vect}[1]{\mathbf{#1}}
\newcommand{\deriv}[2]{\frac{d#1}{d#2}}
\newcommand{\pderiv}[2]{\frac{\partial#1}{\partial#2}}

% ============================================================
% METADATOS DEL DOCUMENTO
% ============================================================
\newcommand{\universidad}{Universidad de Buenos Aires (UBA)}
\newcommand{\facultad}{Facultad de Derecho}
\newcommand{\catedra}{Cátedra de Derechos Reales}
\newcommand{\materia}{Derechos Reales}
\newcommand{\titulodoc}{El Dominio y el Usufructo: de la Teoría a la Práctica Notarial y Societaria}
\newcommand{\autor}{Isaac Edgar Camacho Ocampo}
\newcommand{\anio}{2026}

% ============================================================
% CAJAS ACADÉMICAS
% ============================================================
\newtcolorbox{definition}[1]{
    enhanced,
    breakable,
    title={Definición: #1},
    fonttitle=\bfseries,
    colback=blue!3,
    colframe=blue!50!black,
    boxrule=0.8pt,
    arc=2mm
}

\newtcolorbox{important}{
    enhanced,
    breakable,
    title={Importante},
    fonttitle=\bfseries,
    colback=yellow!8,
    colframe=orange!70!black,
    boxrule=0.8pt,
    arc=2mm
}

\newtcolorbox{example}[1]{
    enhanced,
    breakable,
    title={Ejemplo: #1},
    fonttitle=\bfseries,
    colback=green!4,
    colframe=green!40!black,
    boxrule=0.8pt,
    arc=2mm
}

\newtcolorbox{formula}[1]{
    enhanced,
    breakable,
    title={Fórmula / Regla: #1},
    fonttitle=\bfseries,
    colback=purple!4,
    colframe=purple!50!black,
    boxrule=0.8pt,
    arc=2mm
}

\newtcolorbox{exercise}[1]{
    enhanced,
    breakable,
    title={Ejercicio: #1},
    fonttitle=\bfseries,
    colback=gray!5,
    colframe=gray!60!black,
    boxrule=0.8pt,
    arc=2mm
}

\newtcolorbox{note}{
    enhanced,
    breakable,
    title={Nota},
    fonttitle=\bfseries,
    colback=white,
    colframe=gray!50,
    boxrule=0.6pt,
    arc=2mm
}

% ============================================================
% ENCABEZADOS Y PIES DE PÁGINA
% ============================================================
\pagestyle{fancy}
\fancyhf{}

\fancyhead[L]{\nouppercase{\leftmark}}
\fancyhead[R]{\materia}

\fancyfoot[C]{\thepage}

\renewcommand{\headrulewidth}{0.4pt}
\renewcommand{\footrulewidth}{0pt}

\fancypagestyle{plain}{
    \fancyhf{}
    \fancyfoot[C]{\thepage}
    \renewcommand{\headrulewidth}{0pt}
}

% ============================================================
% HIPERVÍNCULOS Y METADATOS PDF
% ============================================================
\usepackage[
    colorlinks=true,
    linkcolor=blue!50!black,
    citecolor=green!40!black,
    urlcolor=blue!60!black,
    pdftitle={\titulodoc},
    pdfauthor={\autor},
    pdfsubject={Apuntes universitarios},
    bookmarks=true
]{hyperref}

% ============================================================
% DOCUMENTO
% ============================================================
\begin{document}

% ---------------- PORTADA ----------------
\begin{titlepage}
\thispagestyle{empty}

\begin{center}

\vspace*{1cm}

{\large \textsc{\universidad}}

\vspace{0.2cm}

{\large \textsc{\facultad}}

\vspace{0.2cm}

{\large \catedra}

\vspace{3cm}

{\Huge \textbf{\materia}}

\vspace{1cm}

{\LARGE \titulodoc}

\vspace{0.8cm}

\textit{Comprender la lógica del usufructo es aprender a leer, detrás de cada bien, las distintas voluntades que pueden coexistir sobre él.}

\vspace{1.2cm}

\rule{0.70\textwidth}{0.5pt}

\vspace{0.5cm}

{\large Apuntes de estudio}

\vspace{0.5cm}

\rule{0.70\textwidth}{0.5pt}

\vfill

{\large \autor}

\vspace{0.3cm}

{\large \anio}

\end{center}

\vfill

\begin{flushleft}
\textit{Este es un modesto aporte a los estudiantes de ``Derechos Reales'' no pretende sustituir el material oficial de la materia.}
\end{flushleft}

\end{titlepage}

% ---------------- ÍNDICE ----------------
\tableofcontents

% ============================================================
% PRESENTACIÓN
% ============================================================
\chapter*{Presentación}
\addcontentsline{toc}{chapter}{Presentación}

Estos apuntes reconstruyen, en forma ordenada, el desarrollo de una clase universitaria de Derechos Reales dedicada al estudio del dominio y, en particular, del usufructo como derecho real sobre cosa ajena. El desarrollo avanza desde el repaso de las distintas formas que puede asumir el dominio ---pleno, revocable y desmembrado--- hasta el análisis pormenorizado del usufructo: su constitución, sus efectos entre cónyuges, su aplicación en la planificación patrimonial y sucesoria, su uso en el ámbito societario (usufructo de acciones y cuotas sociales) y su extinción, con especial atención a las implicancias impositivas en la Provincia de Buenos Aires.

A lo largo del desarrollo, la exposición conceptual se combina con casos prácticos extraídos de la experiencia notarial y registral. Se destacan en recuadros las referencias normativas y los criterios prácticos expuestos, y se sistematiza el contenido en un cuadro de repaso final bajo el método Cornell.

\section*{Utilidad profesional}
\addcontentsline{toc}{section}{Utilidad profesional}

El usufructo es, junto con la hipoteca y la servidumbre, uno de los derechos reales de mayor aplicación práctica y, sin embargo, uno de los más subutilizados por los profesionales del derecho. Dominar su lógica ---quién tiene el uso y goce, quién conserva la disposición, cómo se constituye, cómo se transmite y cómo se extingue--- permite resolver problemas que exceden el aula: planificar la sucesión de un patrimonio familiar sin perder el control sobre los bienes en vida, estructurar la sucesión de una empresa familiar sin que el fundador pierda el manejo de las decisiones estratégicas, diseñar operaciones inmobiliarias que benefician tanto a quien vende como a quien compra, y asesorar correctamente sobre las consecuencias impositivas de cada una de esas decisiones.

El dominio y el usufructo se entrelazan con el derecho de familia (bienes gananciales, asentimiento conyugal), con el derecho societario (usufructo de acciones, derechos políticos y económicos), con el derecho sucesorio (planificación patrimonial, donaciones con reserva de usufructo) y con el derecho tributario (impuesto a la transmisión gratuita de bienes).

\section*{Utilidad personal}
\addcontentsline{toc}{section}{Utilidad personal}

Cualquier persona que done, herede, compre o venda un inmueble, que tenga bienes gananciales, que forme parte de una sociedad familiar o que quiera proteger su vivienda en la vejez sin perder autonomía económica, puede encontrarse frente a las figuras que se explican en este material. Entender la reserva de usufructo permite, por ejemplo, donar un inmueble a los hijos sin quedar desprotegido; entender el asentimiento conyugal evita firmar (o dejar de firmar) un acto que después puede anularse; y entender la valuación fiscal del usufructo evita pagar de más ---o exponerse a un reclamo--- en una operación inmobiliaria.

\section*{Organización del material}
\addcontentsline{toc}{section}{Organización del material}

El contenido se organiza en cuatro bloques temáticos que avanzan de lo general a lo particular: primero se repasan las formas que puede asumir el dominio; luego se estudian los fundamentos del usufructo como derecho real autónomo; después se exploran sus aplicaciones prácticas en el ámbito patrimonial, familiar y societario; y finalmente se abordan su extinción y sus consecuencias impositivas.

Pensemos en una familia con un inmueble y una empresa familiar. El primer bloque explica qué tipo de dominio tienen sobre esos bienes. El segundo explica cómo esa familia puede desdoblar ese dominio ---por ejemplo, cuando los padres donan la nuda propiedad de la casa a los hijos pero se reservan el derecho de seguir viviendo en ella--- y qué reglas de familia entran en juego (bienes gananciales, asentimiento del cónyuge). El tercero muestra cómo esa misma lógica del usufructo se traslada a la empresa: los padres pueden donar las acciones a sus hijos y reservarse el usufructo, cobrando dividendos y controlando decisiones clave mientras preparan la sucesión del negocio. Y el cuarto cierra el círculo explicando qué ocurre el día en que esa familia decide vender el inmueble o cancelar el usufructo, y cuánto le cuesta impositivamente cada camino.

% ============================================================
% BLOQUE 1
% ============================================================
\chapter{El dominio y sus formas}

\section{Derecho e interdisciplinariedad tecnológica}

El desarrollo de la clase se introdujo con una reflexión sobre el vínculo creciente entre el derecho y la tecnología, a partir de la participación reciente del docente en una jornada sobre inteligencia artificial organizada por un laboratorio universitario especializado en la materia.

El docente sostuvo que \textit{«gran parte de nuestras ciencias hoy están interrelacionadas con la inteligencia artificial»}, del mismo modo en que ya ocurre con la medicina. Señaló que esto no implica únicamente el riesgo de que los profesionales del derecho sean reemplazados por sistemas de inteligencia artificial, sino también ---y sobre todo--- la posibilidad de usar esas herramientas como colaboradoras del ejercicio profesional.

Como ejemplo de un debate ya instalado en la doctrina, se mencionó la discusión sobre si las personas jurídicas pueden integrar el directorio de una sociedad. Se afirmó que, aunque persisten posturas apegadas a conceptos antiguos, \textit{«en algún momento va a llegar que una persona jurídica pueda llegar a formar parte de, o sea, pueda ser directora»}. En la misma línea, se adelantó que los próximos interrogantes del derecho societario girarán en torno a las DAO (organizaciones autónomas descentralizadas) y a las sociedades automatizadas, temas que hoy carecen de regulación específica.

\begin{example}{La inscripción registral y el trabajo interdisciplinario}
Se mencionó una resolución del Registro de la Propiedad Inmueble que permite inscribir adjudicaciones por disolución de la sociedad conyugal instrumentadas por acta notarial, en las que solamente uno de los adjudicatarios solicita la adjudicación. Se calificó esta resolución como una solución práctica y ágil para un problema cotidiano de la práctica notarial.
\end{example}

\begin{important}
\textbf{Referencia normativa ---  Resolución del Registro de la Propiedad Inmueble.} Se cita una resolución del RPI que habilita inscribir adjudicaciones por disolución de la sociedad conyugal, instrumentadas por acta notarial, con la sola solicitud de uno de los adjudicatarios.
\end{important}

De esta introducción se desprende una idea central: el derecho ya no se ejerce en compartimentos estancos. \textit{«El derecho no trabaja con abogados solamente. Hoy el derecho trabaja también con ingenieros y con gente de la ciencia dura»}, en referencia a los especialistas en inteligencia artificial. Se advirtió que, si la elaboración de una ley de sociedades queda en manos de un legislador sin formación en derecho societario, el resultado probablemente no captará los desarrollos doctrinarios y jurisprudenciales acumulados en la materia, de allí la importancia de trabajar en forma interdisciplinaria.

\section{Formas del dominio: pleno, revocable y desmembrado}

El dominio, lejos de ser una figura única y absoluta, admite distintas modalidades. Se repasaron dos categorías de dominio imperfecto ---es decir, no pleno---:

\begin{itemize}
    \item \textbf{Dominio revocable}: aquel sujeto a una condición o a un plazo resolutorio.
    \item \textbf{Dominio desmembrado}: aquel en el que alguna de las facultades que integran el dominio ---uso, goce y disposición--- se transfiere a otra persona.
\end{itemize}

\section{El dominio desmembrado y los derechos reales sobre cosa ajena}

\begin{definition}{Dominio desmembrado}
El dominio se integra por tres facultades: uso, goce y disposición. Existe desmembramiento cuando alguna de esas facultades se traslada a otro sujeto, lo cual puede ocurrir mediante la constitución de otro derecho real, aunque no siempre el desmembramiento implica necesariamente esa constitución.
\end{definition}

El derecho real que permite ese desdoblamiento es el usufructo, que otorga a su titular \textit{«el uso y el goce»}. A partir de esta idea se planteó el eje central del desarrollo: el usufructo es \textit{«un derecho real superimportante que la gente no lo sabe aprovechar»}, ya que suele asociarse únicamente al esquema familiar clásico, sin advertir que admite otros usos de gran interés práctico.

Si el dominio se integra por uso, goce y disposición, y el usufructo otorga solamente uso y goce, la disposición permanece en cabeza del titular de dominio, a quien se identifica con el nombre técnico de \textbf{nudo propietario}.

% ============================================================
% BLOQUE 2
% ============================================================
\chapter{El usufructo: fundamentos}

\section{Concepto de usufructo: nuda propiedad y reserva}

Se presentó el caso paradigmático: aquel en el que una persona es única titular de un inmueble y transmite la nuda propiedad reservándose el usufructo para sí.

\begin{definition}{Reserva de usufructo}
Cuando el único titular de un inmueble transfiere o transmite la nuda propiedad, y se reserva para sí el usufructo, se configura la denominada \textbf{reserva de usufructo}: la titularidad de dominio (la propiedad ``nuda'', vacía de uso) se transmite a otra persona, mientras que el uso y goce permanece, por reserva, en cabeza de quien transmite.
\end{definition}

\begin{example}{El lápiz}
Para ilustrar el desdoblamiento entre nuda propiedad y usufructo se utilizó un objeto simple: si una persona tiene un lápiz y le da a otra la nuda propiedad, se queda para sí con el uso del lápiz. El dueño del lápiz pasa a ser el receptor de la nuda propiedad; quien transmite se queda con el uso.
\end{example}

\section{Bienes gananciales y titularidad registral}

Se planteó una pregunta central: ¿puede una persona casada, titular registral de un bien ganancial, reservarse el usufructo para sí y, además, constituirlo a favor de su cónyuge?

La respuesta desarrollada fue negativa en cuanto a la reserva a favor del cónyuge: \textit{«No me puedo reservar. Porque yo me puedo reservar para mí, pero no le puedo reservar a mi cónyuge, porque mi cónyuge es titular de dominio.»}

\begin{important}
Aunque un bien sea ganancial, la \textbf{titularidad registral} corresponde exclusivamente a quien figura como titular en la escritura. La gananciabilidad solamente produce efectos patrimoniales (participación en el 50\%) en dos momentos: al momento de la disolución de la sociedad conyugal o al momento del fallecimiento. Mientras esas dos situaciones no se produzcan, la administración del bien corresponde en forma exclusiva al titular registral.
\end{important}

Un matrimonio puede adquirir un bien durante la vigencia de la sociedad conyugal, pero si la escritura se otorga únicamente a nombre de uno de los cónyuges, no es posible subsanar esa situación mediante un contrato entre cónyuges para cambiar la titularidad: quien registró el bien a su nombre conserva la administración durante todo el matrimonio. Como se sintetizó en el desarrollo: \textit{«En nuestra cabeza está que si es ganancial, es de los dos. Sí, es de los dos desde lo patrimonial, pero no desde la realidad fáctica. ¿Por qué? Porque la administración la tiene solamente el titular, y solamente va a recaer sobre los dos cuando se parte.»}

% TODO: contenido incompleto o ambiguo en las notas originales — no se detalla en el apunte con qué frecuencia ni por qué motivos concretos las parejas separadas o divorciadas mantienen sin disolver la sociedad conyugal, más allá de la mención genérica a razones económicas o de otros vínculos societarios.

\section{Constitución de usufructo entre cónyuges}

Descartada la reserva directa a favor del cónyuge, se explicaron dos alternativas para lograr un resultado equivalente.

La primera vía es el cambio del régimen patrimonial matrimonial: puede optarse por mutar del régimen de comunidad de gananciales al régimen de separación de bienes, lo que habilitaría una adjudicación del 50\% a cada cónyuge. Sin embargo, se relativizó la utilidad práctica de esta vía, dado que la mayoría de las personas no tiene interés en modificar todo su régimen patrimonial matrimonial solo para otorgarle el usufructo al cónyuge.

\begin{formula}{Doble operación simultánea entre cónyuges}
La segunda vía consiste en una doble operación simultánea: el titular registral transmite el dominio a la otra persona y, en el mismo acto, la otra persona le constituye el usufructo a quien transmitió. Aplicado a una cuota parte: el titular puede reservarse el usufructo sobre su 50\% y, en el mismo acto, el nuevo nudo propietario constituye el usufructo sobre el otro 50\% a favor de su cónyuge.

El resultado práctico es que ambos cónyuges terminan siendo cousufructuarios del inmueble en las proporciones acordadas, aun cuando uno solo de ellos conserve la nuda propiedad plena. Esta estructura resulta particularmente útil en operaciones con montos elevados de impuesto de sellos, donde la ingeniería de la operación puede incidir en la carga fiscal.
\end{formula}

\section{El asentimiento conyugal}

Se introdujo un supuesto complementario: la constitución de un usufructo a favor del cónyuge del nudo propietario, cuando quien dona no tiene, al momento de donar, cónyuge, pero lo adquiere con posterioridad (por ejemplo, una donación a un hijo casado, sobre la que luego se pretende constituir un usufructo a favor de la nuera o el yerno).

Si el hijo (nudo propietario) está casado, y el bien fue recibido por donación (es decir, es un bien propio para él), la constitución de un derecho real sobre ese bien requiere, en principio, la conformidad del cónyuge, porque se está disponiendo del bien, salvo que se tratara de un bien propio con conformidad ya prestada por la cónyuge.

\begin{important}
\textbf{Criterio ---  Asentimiento conyugal y bienes propios.} Cuando se constituye, transfiere o modifica un derecho real, se requiere el asentimiento conyugal, con la excepción de los bienes propios que no constituyan la sede del hogar conyugal. Si el bien propio es la sede del hogar conyugal, el asentimiento se exige igualmente, con el fin de proteger la vivienda familiar.
\end{important}

En el propio acto de donación, si el bien resulta propio para el donatario por provenir de una liberalidad, suele requerirse igualmente que el cónyuge del donatario ratifique o reconozca ese carácter propio.

\begin{example}{Observación registral sobre una partición sucesoria}
Una escritura instrumentaba una adjudicación por partición derivada de una sucesión, sin que compareciera el o la cónyuge del adjudicatario a prestar el asentimiento conyugal. El escribano interviniente no exigió ese asentimiento por entender que no se trataba de la sede del hogar conyugal y que el carácter propio del bien resultaba indudable, por provenir de una sucesión.

El acto fue observado por el organismo registral, que exigió el asentimiento conyugal. Se manifestó desacuerdo explícito con esa observación: \textit{«A mi criterio, estaba mal observado, porque no era sede del hogar conyugal y era clarísimo que era un bien propio, porque provenía justamente de una sucesión, partición, adjudicación.»}

Sobre el criterio aplicado en la práctica profesional: \textit{«De hecho, yo no pido el asentimiento conyugal cuando viene de bienes propios, porque es claro. Ahora, si dudo dónde viven ---si coincide el domicilio del acto con el domicilio del inmueble que se está enajenando--- ahí sí le pido el asentimiento conyugal.»}
\end{example}

De este modo, la exigencia del asentimiento conyugal cede frente a la certeza sobre el carácter propio del bien y sobre el hecho de que no constituye la sede del hogar conyugal; ante la duda sobre este último punto corresponde requerir igualmente la conformidad del cónyuge, como resguardo de la vivienda familiar.

% ============================================================
% BLOQUE 3
% ============================================================
\chapter{El usufructo en la práctica patrimonial y societaria}

\section{El usufructo oneroso}

Se distinguieron dos momentos ya vistos ---cómo se adquiere el usufructo por reserva o por constitución--- y una tercera vía: la compra del usufructo. Se introdujo el tema señalando que es allí donde se advierte la riqueza de este derecho real, subutilizado en la práctica.

\begin{example}{El propietario de 85 años y el inquilino de toda la vida}
Un propietario de 85 años que alquila desde hace diez años un local o departamento al mismo inquilino le ofrece venderle la nuda propiedad: para cuando el propietario fallezca, el inmueble queda para el inquilino. El propietario se reserva el usufructo, sigue cobrando el alquiler como siempre, pero le da al inquilino la tranquilidad de que el día en que fallezca, el inmueble será suyo. El precio de venta se reduce respecto del valor pleno del inmueble.

La lógica económica de esa quita sobre el valor responde a que el precio de la nuda propiedad se reduce en función de la expectativa de vida del usufructuario, funcionando de modo análogo a una renta vitalicia, pero cobrada en un pago único en lugar de en cuotas periódicas.
\end{example}

\begin{example}{La inversión de un gran inversor sobre un terreno}
Un nudo propietario recibe una oferta de un inversor de gran escala interesado en desarrollar una inversión de magnitud sobre un terreno. En lugar de un contrato de locación por diez años ---que no lo resguardaría frente al fallecimiento del propietario o a un eventual embargo---, el inversor le compra el usufructo del terreno a título oneroso.

De este modo, el inversor adquiere por compraventa el usufructo del terreno, blindando su inversión frente a contingencias sucesorias o de embargos que podrían afectar a un simple contrato de locación.
\end{example}

\begin{formula}{Plazo del usufructo oneroso}
El usufructo adquirido a título oneroso puede pactarse hasta el momento del fallecimiento del usufructuario, o con un plazo cierto y determinado. Si el usufructuario es una persona humana, el usufructo se extingue con su fallecimiento, aun cuando se hubiera pactado un plazo fijo mayor.

En consecuencia, si se prevé adquirir un usufructo, conviene que el adquirente sea una persona jurídica con un plazo de duración mayor al tiempo pactado del usufructo: en tal caso, la extinción se producirá por cancelación de esa sociedad, y no por fallecimiento.
\end{formula}

\section{Extinción y derecho de acrecer entre usufructuarios}

El usufructo \textit{«se extingue con el fallecimiento del usufructuario. No es que puede haber sucesión, no es que se lo pueda testar, no es que se van a incorporar mis herederos.»}

El punto se complejiza cuando hay pluralidad de usufructuarios, a partir de la existencia o no de derecho de acrecer entre los cousufructuarios.

\begin{definition}{Derecho de acrecer entre cousufructuarios}
Si varios usufructuarios \textbf{no} tienen entre sí derecho de acrecer, el usufructo se va extinguiendo parcialmente a medida que cada uno fallece, y esa porción consolida el dominio pleno en cabeza del nudo propietario. Si \textbf{sí} se pactó derecho de acrecer, el o los usufructuarios sobrevivientes se quedan con la totalidad del usufructo.
\end{definition}

\begin{important}
Es frecuente omitir, por error, la inclusión del derecho de acrecer al constituir un usufructo a favor de varios cousufructuarios (por ejemplo, ambos cónyuges donantes que se reservan el usufructo). Si no se incluye el derecho de acrecer y fallece uno de los usufructuarios, el nudo propietario adquiere el dominio pleno sobre la parte ideal correspondiente a ese usufructuario. Los cousufructuarios no necesariamente deben ser cónyuges: pueden ser convivientes o incluso personas sin ningún vínculo afectivo entre sí, por lo que la inclusión (o no) del derecho de acrecer debe decidirse expresamente en cada caso.
\end{important}

\begin{example}{Acrecimiento proporcional entre tres usufructuarios}
Tres personas son cousufructuarias de un mismo bien: la primera titular del 50\% del usufructo, y las otras dos titulares del 25\% cada una. Al fallecer la primera usufructuaria, si existe derecho de acrecer, su 50\% no se reparte en partes iguales entre las dos restantes, sino a prorrata de la proporción que cada una ya tenía sobre el 50\% remanente. Así, si la segunda usufructuaria tenía el 30\% de ese 50\% restante y la tercera el 20\%, cada una acrece en esa misma proporción sobre el porcentaje liberado, y no en partes iguales.
\end{example}

\section{Usufructo de acciones y cuotas sociales}

El análisis se trasladó del inmueble a otro tipo de bien: las acciones o cuotas sociales de una sociedad. En lo estructural, el usufructo de acciones sigue reglas equivalentes a las del usufructo inmobiliario: si la acción es un bien ganancial, corresponde el asentimiento conyugal en los mismos supuestos y con las mismas excepciones ya analizadas; y el titular puede reservarse el usufructo o constituirlo a favor de un tercero, sin distinciones respecto del régimen aplicable al dominio inmobiliario.

\begin{definition}{Derechos políticos y derechos económicos}
La acción de una sociedad comprende dos categorías de derechos: los \textbf{derechos económicos} (dividendos, cuota de liquidación) y los \textbf{derechos políticos} (voto en las asambleas). Históricamente, la doctrina asociaba el usufructo de acciones exclusivamente a los derechos económicos, dejando los derechos políticos en cabeza del nudo propietario.
\end{definition}

\begin{important}
\textbf{Referencia normativa ---  Ley de Sociedades.} Se hace referencia a la Ley 19.550 de Sociedades Comerciales y a los proyectos de reforma societaria en debate, que regulan de manera distinta la posición del usufructuario minoritario frente a las decisiones políticas de la sociedad.
\end{important}

\begin{example}{Donación de acciones sin reserva de derechos políticos}
Un padre dona a su hijo las acciones de la empresa familiar, con fines de planificación sucesoria, reservándose únicamente los derechos económicos (dividendos y cuota de liquidación), sin advertir que estaba cediendo también los derechos políticos.

Como consecuencia, el hijo ---titular de los derechos políticos--- puede votar en una asamblea la disolución y liquidación de la sociedad, con lo cual el padre queda con derecho a una cuota de liquidación, pero pierde la posibilidad de seguir cobrando dividendos, que era lo que efectivamente pretendía conservar: \textit{«Yo no quiero una cuota de liquidación: yo quiero que la sociedad permanezca, quiero seguir cobrando dividendos.»}

Otros supuestos en los que el nudo propietario, ejerciendo en soledad los derechos políticos, puede licuar o extinguir la expectativa económica del usufructuario: un aumento de capital que diluye su porcentaje de participación (por ejemplo, del 50\% original a un 10\% tras el aumento decidido por los demás socios), una fusión societaria frente a la cual el usufructuario solo puede ejercer el derecho de receso, o directamente el cambio del objeto social hacia una actividad distinta de la original.
\end{example}

La conclusión doctrinaria que se sostuvo sobre el punto fue la siguiente: \textit{«Cuando se trata de derechos políticos, no es una caja vacía de contenido cuando yo te doy el usufructo. [...] Yo me quedé con el uso y el goce, vos vas a tener la nuda propiedad, pero yo ya te estoy dando algo como nudo propietario.»}

Se explicó, en la misma línea, por qué existe un mercado real para la compraventa exclusiva de la nuda propiedad, sin el usufructo: quien compra una nuda propiedad con un usufructuario de noventa años a un precio muy reducido está haciendo una apuesta financiera razonable sobre el tiempo de duración esperado del usufructo. Esta práctica se señaló como habitual en España, comparándosela con la hipoteca inversa.

\begin{note}
Como alternativa superadora al esquema de todo o nada, es posible graduar en el tiempo la cesión de derechos económicos y políticos, en función de la maduración del sucesor dentro de la empresa: se dona la acción reservándose el usufructo de los derechos políticos y económicos, y progresivamente, a medida que transcurren los años, se cede un porcentaje creciente de esos derechos económicos reservados. La justificación de esta gradualidad es formativa: permite que quien recibe la nuda propiedad se prepare para ejercer, en el futuro, los derechos políticos en la asamblea.
\end{note}

Sobre la comparecencia en la asamblea cuando existe usufructo sobre las acciones, la regla depende de a quién se le reservaron los derechos políticos: si el usufructuario conserva esos derechos, es quien comparece y vota; el nudo propietario, en tal caso, no comparece, y solo percibe los derechos económicos que le correspondan. Cuando existe comparecencia conjunta en condominio de derechos, rige la regla de unificación de la voluntad de voto: \textit{«Se tiene que unificar la votación en uno, porque una acción, un voto.»}

\begin{definition}{Fideicomiso accionario}
Estructura mediante la cual las acciones ---gravadas con usufructo y nuda propiedad--- se aportan a un fideicomiso, cuyo fiduciario ejerce los derechos políticos conforme a las pautas fijadas en el instrumento de constitución, en lugar de que cada cousufructuario o nudo propietario comparezca individualmente a la asamblea.
\end{definition}

\section{Planificación sucesoria de la empresa familiar}

El usufructo de acciones constituye una herramienta de planificación sucesoria empresarial que permite al fundador de una empresa familiar transmitir la propiedad de las acciones a la siguiente generación ---con el consiguiente beneficio de ordenar el patrimonio en vida y evitar conflictos futuros entre herederos--- sin resignar, necesariamente, el control económico ni político del negocio mientras viva y lo considere necesario.

La combinación de reserva de usufructo, derechos políticos condicionados a determinados actos societarios (fusión, aumento de capital, cambio de objeto social) y cesión progresiva de derechos económicos conforma un esquema más rico y flexible que la simple donación lisa y llana de las acciones.

% ============================================================
% BLOQUE 4
% ============================================================
\chapter{Aspectos fiscales y extinción del usufructo}

\section{Cancelación del usufructo por renuncia}

Estando con vida el usufructuario, el usufructo se cancela mediante la renuncia del usufructuario a su derecho, lo que consolida el dominio pleno en cabeza del nudo propietario: \textit{«Yo renuncio a ese derecho real que tengo a mi favor; le queda el dominio pleno a favor de mi hijo. Mi hijo vende él el departamento.»}

\begin{important}
\textbf{Referencia normativa ---  Impuesto a la Transmisión Gratuita de Bienes (Provincia de Buenos Aires).} La Provincia de Buenos Aires es, dentro del país, la única jurisdicción que mantiene vigente el impuesto a la transmisión gratuita de bienes, aplicable a los actos a título gratuito. Las causales que determinan su procedencia son la ubicación del bien en la Provincia o el domicilio del causante o del sucesor en dicha jurisdicción.
\end{important}

\begin{example}{Renuncia al usufructo: Ciudad de Buenos Aires frente a Provincia de Buenos Aires}
Una madre usufructuaria y su hijo nudo propietario deciden vender un inmueble.

Si el inmueble está ubicado en la Ciudad de Buenos Aires, la renuncia al usufructo no genera preocupación tributaria adicional, porque esa jurisdicción no aplica el impuesto a la transmisión gratuita de bienes: la madre renuncia gratuitamente a su derecho, el hijo queda como titular del dominio pleno y puede vender libremente el inmueble.

En cambio, en la Provincia de Buenos Aires, la renuncia gratuita al usufructo configura, a criterio fiscal, un acto a título gratuito que beneficia patrimonialmente al nudo propietario, y tributa el impuesto a la transmisión gratuita de bienes.
\end{example}

\begin{formula}{Alternativa a la renuncia gratuita en Provincia de Buenos Aires}
En lugar de que la usufructuaria renuncie gratuitamente y luego el hijo venda el dominio pleno, puede estructurarse una doble venta a título oneroso al mismo comprador: el nudo propietario le vende la nuda propiedad, y el usufructuario le vende, en el mismo acto, el usufructo. De este modo no hay renuncia gratuita a la valuación, y el efecto económico final es equivalente (el comprador termina con el dominio pleno) pero sin que medie el acto gratuito que dispara el impuesto provincial.

En la Ciudad de Buenos Aires, a los fines de los honorarios notariales, la venta del usufructo tributa el 1\% y la venta de la nuda propiedad tributa el 2,7\%.
\end{formula}

\begin{important}
El uso de la diferencia de alícuotas entre la venta del usufructo y la venta de la nuda propiedad debe manejarse con cuidado: recurrir a esta estructura únicamente por motivos impositivos, para pagar menos impuesto, requiere especial prudencia profesional.
\end{important}

\section{Valuación del usufructo a efectos impositivos}

El cierre del desarrollo se dedicó a un problema práctico de relevancia notarial: cómo se distribuye el precio total de una operación entre el valor asignado al usufructo y el valor asignado a la nuda propiedad, cuando ambos se venden por separado al mismo comprador.

\begin{example}{Subvaluación indebida del usufructo}
Algunos profesionales, buscando reducir la carga del impuesto de sellos (más bajo sobre la cuota del usufructo), asignan al usufructo un valor artificialmente bajo: por ejemplo, vender el usufructo en diez mil y la nuda propiedad en cincuenta mil, sin relación con sus valores reales. Esta práctica no debe aceptarse.

El fundamento del rechazo es técnico: el valor del usufructo debe guardar proporción con la expectativa de vida ---y, por lo tanto, de uso y goce efectivo--- del usufructuario; no puede fijarse arbitrariamente. Quien perjudica esta práctica de subvaluación es, fundamentalmente, el comprador, que termina pagando por la nuda propiedad un precio de sellos superior al que en rigor le correspondería si la proporción entre usufructo y nuda propiedad reflejara razonablemente la expectativa de vida del usufructuario.
\end{example}

\begin{formula}{Proporción máxima orientativa del usufructo}
Como criterio orientativo y sujeto a la edad del usufructuario, el valor asignado al usufructo no debería superar, como máximo, entre el 25\% y el 30\% del valor total de la operación. La proporción se calcula sobre el valor total de la venta, y no sobre una porción aislada de la operación.
\end{formula}

\begin{exercise}{Distribución del precio en una operación con usufructo}
Sobre una operación de \$200.000, con una alícuota del impuesto de sellos del 2,7\%, se planteó cómo distribuir el precio entre el usufructo y la nuda propiedad. Frente a la propuesta de calcular el 30\% sobre la porción del usufructo aisladamente, se corrigió el criterio: la proporción se calcula sobre el total de la operación, y no sobre una parte separada, ya que ambas ventas ---usufructo y nuda propiedad--- integran una misma operación económica y deben evaluarse conjuntamente a los fines de determinar si la distribución del precio resulta razonable.
% TODO: contenido incompleto o ambiguo en las notas originales — el apunte no consigna el resultado numérico final del ejercicio, solo el criterio de cálculo corregido en clase.
\end{exercise}

El riesgo central de estas prácticas de subvaluación se resumió del siguiente modo: en el manejo impositivo se pueden generar situaciones que afectan al usufructuario y al nudo propietario, pero sobre todo al comprador, que resulta en general el más perjudicado.

% ============================================================
% CUADRO CORNELL DE REPASO
% ============================================================
\chapter{Cuadro Cornell de repaso}

El siguiente cuadro sintetiza, bajo el formato de notas Cornell, las preguntas y conceptos clave desarrollados a lo largo del material, junto con su respuesta o explicación, cerrando con una síntesis integradora de todo el contenido.

\begin{longtable}{
    >{\raggedright\arraybackslash}p{0.30\textwidth}
    >{\raggedright\arraybackslash}p{0.60\textwidth}
}
\toprule
\textbf{Preguntas / palabras clave} & \textbf{Notas / respuestas} \\
\midrule
\endfirsthead

\toprule
\textbf{Preguntas / palabras clave} & \textbf{Notas / respuestas} \\
\midrule
\endhead

\midrule
\multicolumn{2}{r}{\textit{Continúa en la página siguiente}} \\
\endfoot

\bottomrule
\endlastfoot

\textbf{¿Qué formas puede asumir el dominio además del dominio pleno?} &
El dominio puede ser revocable (sujeto a condición o plazo resolutorio) o desmembrado (cuando alguna de las facultades de uso, goce o disposición se transfiere a otra persona, generalmente mediante la constitución de un derecho real sobre cosa ajena, como el usufructo). \\

\textbf{¿Qué es la nuda propiedad y qué es la reserva de usufructo?} &
La nuda propiedad es la titularidad de dominio despojada del uso y goce. La reserva de usufructo ocurre cuando el único titular de un bien lo transmite reteniendo para sí el uso y goce, mientras el adquirente recibe la nuda propiedad. \\

\textbf{¿Puede una persona reservarse el usufructo a favor de su cónyuge?} &
No. Cada cónyuge solo puede reservarse el usufructo para sí mismo, porque el otro cónyuge no es titular de dominio del bien (aun siendo ganancial). Para que el cónyuge acceda al usufructo, el titular debe transmitir la nuda propiedad y, en el mismo acto, el nuevo nudo propietario debe constituir el usufructo a su favor. \\

\textbf{Un bien es ganancial pero está a nombre de un solo cónyuge: ¿quién lo administra?} &
Lo administra únicamente el titular registral. La calidad de ganancial solo produce efectos patrimoniales (participación en el 50\%) al momento de la disolución de la sociedad conyugal o del fallecimiento; hasta entonces, la administración es exclusiva del titular. \\

\textbf{¿Cuándo se requiere el asentimiento conyugal para constituir un derecho real?} &
Se requiere siempre que se disponga de un bien ganancial. Respecto de bienes propios, no se requiere, salvo que el bien constituya la sede del hogar conyugal, en cuyo caso el asentimiento se exige igualmente para proteger la vivienda familiar. \\

\textbf{¿Cómo se adquiere un usufructo a título oneroso?} &
Mediante compraventa: el nudo propietario le vende el usufructo a un tercero interesado (por ejemplo, un inversor), o el usufructuario le vende la nuda propiedad al tenedor del inmueble (por ejemplo, a un inquilino de largo plazo). El precio se ajusta según la expectativa de vida del usufructuario. \\

\textbf{¿Por qué conviene que el usufructo oneroso lo adquiera una persona jurídica y no una persona humana?} &
Porque el usufructo se extingue con el fallecimiento del usufructuario si es persona humana, aun cuando se hubiera pactado un plazo mayor. Si el usufructuario es una persona jurídica, la extinción se produce recién con la cancelación de esa sociedad, permitiendo plazos de duración más largos y previsibles. \\

\textbf{¿Qué es el derecho de acrecer entre cousufructuarios?} &
Es la cláusula que determina qué ocurre al fallecer uno de varios usufructuarios: si no hay derecho de acrecer, esa porción consolida el dominio pleno en el nudo propietario; si lo hay, los usufructuarios sobrevivientes se quedan con la totalidad (o la porción liberada) a prorrata de su participación previa. \\

\textbf{¿Qué diferencia hay entre derechos políticos y derechos económicos en el usufructo de acciones?} &
Los derechos económicos (dividendos, cuota de liquidación) suelen quedar en cabeza del usufructuario; los derechos políticos (voto en asambleas) pueden reservarse total o parcialmente, ya que si quedan en manos del nudo propietario sin control, este puede votar decisiones (fusión, aumento de capital, cambio de objeto) que licúen o extingan la expectativa económica del usufructuario. \\

\textbf{¿Cómo se puede planificar progresivamente la sucesión de una empresa familiar mediante el usufructo?} &
El fundador dona las acciones reservándose el usufructo (derechos políticos y económicos) y, con el paso de los años, va cediendo gradualmente porcentajes de derechos económicos y de participación política a medida que el sucesor se forma dentro de la empresa, sin resignar el control mientras lo considere necesario. \\

\textbf{¿Qué es un fideicomiso accionario y para qué se usa en este contexto?} &
Es una estructura mediante la cual las acciones gravadas con usufructo y nuda propiedad se aportan a un fideicomiso, cuyo fiduciario ejerce los derechos políticos según reglas predefinidas, evitando la necesidad de unificar votos entre cousufructuarios o nudos propietarios en cada asamblea. \\

\textbf{¿Cómo se cancela un usufructo en vida del usufructuario?} &
Mediante renuncia expresa del usufructuario a su derecho real, lo que consolida automáticamente el dominio pleno en cabeza del nudo propietario. \\

\textbf{¿Qué es el impuesto a la transmisión gratuita de bienes y dónde rige?} &
Es un impuesto que grava los actos a título gratuito y que hoy solo se mantiene vigente en la Provincia de Buenos Aires. Sus causales de procedencia son la ubicación del bien en la Provincia o el domicilio del causante/sucesor en dicha jurisdicción. \\

\textbf{¿Cómo se evita el impuesto a la transmisión gratuita al vender un inmueble con usufructo en Provincia de Buenos Aires?} &
En lugar de que el usufructuario renuncie gratuitamente al usufructo (acto gravado) y luego el nudo propietario venda el dominio pleno, se estructuran dos ventas onerosas simultáneas al mismo comprador: una del usufructo y otra de la nuda propiedad. \\

\textbf{¿Qué proporción del precio total puede asignarse razonablemente al usufructo en una operación?} &
Como criterio orientativo, no más del 25\% al 30\% del valor total de la operación, calculado siempre sobre el total y no sobre una porción aislada, en función de la expectativa de vida del usufructuario. Subvaluar el usufructo para pagar menos impuesto de sellos perjudica principalmente al comprador. \\

\midrule
\multicolumn{2}{p{0.90\textwidth}}{
\textbf{Resumen:} El desarrollo reconstruye, a partir del repaso de las formas imperfectas del dominio, la lógica completa del usufructo como derecho real de uso y goce sobre cosa ajena: cómo se constituye (por reserva o por constitución expresa), qué límites impone el régimen de bienes del matrimonio (titularidad registral, asentimiento conyugal), cómo se adquiere a título oneroso y por qué conviene resguardar esa adquisición en una persona jurídica, qué ocurre con la pluralidad de usufructuarios y el derecho de acrecer, cómo se traslada toda esta lógica al ámbito societario distinguiendo derechos políticos de derechos económicos ---convirtiéndose en una herramienta central de planificación sucesoria empresarial---, y finalmente cómo se extingue el usufructo y qué impacto impositivo tiene esa extinción, en particular el impuesto a la transmisión gratuita de bienes vigente en la Provincia de Buenos Aires. El usufructo es una herramienta rica y flexible, subutilizada en la práctica profesional, capaz de articular en un mismo instrumento los intereses patrimoniales, familiares, societarios y fiscales de una persona o una familia a lo largo del tiempo, sin necesidad de resignar anticipadamente el control sobre los bienes propios.
} \\

\end{longtable}

\end{document}
